% R2 manuscript: methods-and-evaluation revision of Paper 2 (CALM).
% Decision: mjeragh/VANET#67, 8 November 2026; plan: VANET/MyScenarios/AnyScenario/Documentation/PAPER2_R2_MANUSCRIPT_PLAN.md
% IEEEtran is a neutral carrier until the venue is chosen. The root main.tex is the withdrawn T-IV record and is not edited.
% Tables in tables/ are generated by VANET/.../paper2_revision/make_r2_tables.py from the record files; do not edit them by hand.
\documentclass[journal]{IEEEtran}

\usepackage{amsmath,amssymb,amsfonts,amsthm}
\usepackage{graphicx}
\usepackage{textcomp}
\usepackage{booktabs}
\usepackage{array}
\usepackage{multirow}
\usepackage{xcolor}
\usepackage{hyperref}
\usepackage{cite}
\usepackage{bm}

\newtheorem{theorem}{Theorem}
\newtheorem{definition}{Definition}

\newcommand{\KL}{\mathrm{KL}}
\newcommand{\E}{\mathbb{E}}
\newcommand{\todo}[1]{\textcolor{red}{[#1]}}

\begin{document}

\title{What the Benchmark Hid: A Corrected Evaluation of Imitation-Anchored Multi-Agent Reinforcement Learning for Cooperative Highway Exits}
% Alternative: Anchoring Preserves, Fine-Tuning Does Not Improve: Evaluation Pitfalls in SUMO + NGSIM Studies of Imitation-Anchored MARL

\author{Mohammad~Jeragh~and~Ebrahim~Alrashed%
\thanks{M. Jeragh and E. Alrashed are with the Department of Computer Engineering, College of Engineering and Petroleum, Kuwait University, P.O.~Box 5969 Safat, 13060 Kuwait (e-mail: mohammad.jeragh@ku.edu.kw; e.alrashed@ku.edu.kw).}}

\maketitle

\begin{abstract}
We previously reported that CALM, a multi-agent reinforcement learning (MARL) method that pre-trains a behaviour-cloned (BC) prior on NGSIM driving data and fine-tunes it under a decaying KL anchor, prevents mode collapse and produces diverse, human-like and safe cooperative highway-exit behaviour in SUMO. Auditing that study's pipeline, we found three input/action-contract defects (the prior and the fine-tuned actor consumed different feature spaces, the BC action order was never permuted, and the lane feature was one NGSIM lane off), two evaluator defects (agents were dropped after the first step; SUMO, not the policy, executed the agents' lane changes), an in-sample NGSIM comparison, and a missing majority-class baseline. Together they made a degenerate constant policy look like a diverse, safe and human-like one. We rebuilt the pipeline under a single declared contract (vehicle-grouped NGSIM split with a frozen test set, Definition~1 entropy estimated over states, declared execution rule, policy-attributable lane changes, exposure-normalised proximity, constant-action and simulator-default reference policies) and re-ran the evaluation with pre-registered decision rules. Three things are measurable in this benchmark and were measured. First, fine-tuning the BC prior without the anchor collapses it into a low-entropy lane oscillator (per-state entropy 0.22--0.25, 38--45\,\% exits) on both seeds; with the anchor it stays at 0.71 in the simulator and 0.81 on held-out NGSIM states, and the pre-registered rule for ``anchoring prevents collapse'' passes. Second, on matched episodes the anchored policy is no better than the frozen prior it protects on any declared measure: the prior has equal-or-higher entropy, more exits (99\,\% vs 80--85\,\%) and the same speed; fine-tuning trained out the prior's lane oscillation and gained nothing else. Third, greedy evaluation of the anchored policy is the constant Lane\_Right policy, zero collisions is a property of SUMO's car-following for every policy including stranded ones, and NGSIM action agreement is below the 55\,\% majority baseline for every method. The anchor is a regulariser that preserves, not one that improves, and the benchmark cannot test human-likeness, safety or cooperation. We give the defect catalogue, the corrected contract and a reporting checklist for imitation-plus-RL driving evaluations in SUMO against NGSIM.
\end{abstract}

\begin{IEEEkeywords}
Multi-agent reinforcement learning, behavioural cloning, mode collapse, evaluation methodology, SUMO, NGSIM, reproducibility
\end{IEEEkeywords}

\section{Introduction}
\todo{Plan \S2 row 1. The submitted claims (review \S2 table), what the corrected evaluation found, honest framing as the correction of our own study.}

Contributions:
\begin{enumerate}
  \item A defect catalogue for an imitation-plus-RL driving pipeline, each defect with its measured effect on the reported numbers (Section~\ref{sec:defects}).
  \item Three identity results that any such study must check: greedy evaluation of the anchored policy equals the constant Lane\_Right policy; zero collisions equals SUMO's car-following, for every policy; NGSIM action agreement below the majority baseline for every method (Section~\ref{sec:concealed}).
  \item A pre-registered isolation of the anchoring effect: fixed initialisation, one training and evaluation contract, matched episode seeds, a decision rule written before the runs (Section~\ref{sec:anchor}).
  \item The frozen-prior control on the same matched episodes (Section~\ref{sec:prior}).
  \item A testability checklist for the MultiExit + NGSIM benchmark and a reporting checklist for SUMO + NGSIM evaluations (Section~\ref{sec:testability}).
\end{enumerate}

\section{Background and Related Work}
\label{sec:background}
\todo{CALM as submitted: BC prior, decaying KL anchor, Definition 1, Theorem 1 kept; Proposition 1 removed (plan A8). KL-to-prior fine-tuning: AWAC, BRAC, TD3+BC, KL-regularised RLHF. Evaluation pitfalls in deep RL; IL evaluation leakage; SUMO safety layers.}

% Section 3: the submitted pipeline and its defects.
% Evidence: VANET protocol §1, §4–§4d, §8.7 finding 4, §8.8 findings 1 and 5; ESWA plan §2b–§2d, §9.1–§9.3; scope review §1–§2.
\section{The Submitted Pipeline and Its Defects}
\label{sec:defects}

\subsection{The pipeline as submitted}
\label{sec:pipeline}

The submitted study trained five MADDPG actors for the MultiExit scenario in SUMO~\cite{lopez2018microscopic}, a four-lane highway segment with three exit ramps in which three of the five agents are assigned an exit and two are mainline vehicles. Each actor was first fitted by behavioural cloning to NGSIM Interstate~80 trajectories~\cite{ngsim2006} discretised into five actions (Lane\_Left, Lane\_Right, Accelerate, Decelerate, Maintain), then fine-tuned by MADDPG under a KL penalty towards a frozen copy of the cloned network, with a weight $\lambda(t)=\lambda_0\delta^t$ that decays over episodes. Two evaluators produced the reported numbers: a \emph{NGSIM validator} that replays the policy on recorded human states and reports action agreement, Jensen--Shannon similarity and an entropy, and a \emph{safety evaluator} that rolls the policy out in SUMO and reports exit success, collisions, average speed and close-proximity events. Training reward was taken from the training log.

Auditing this pipeline against its own git history and run logs, we found that the three stages---imitation prior, fine-tuning and evaluation---did not share an input or action contract, that both evaluators measured something other than what the paper said, and that the comparison against NGSIM was in-sample. Table~\ref{tab:t1_defects} lists the defects with the mechanism, the measured effect on the submitted numbers and the correction described in Section~\ref{sec:contract}. The rest of this section groups them.

\subsection{Input and action contract (D1--D3)}
\label{sec:defects_contract}

\paragraph{Feature spaces (D1).}
The cloned prior was fitted to four features computed from NGSIM and normalised to $[0,1]$: lane position, velocity, distance to the exit and headway. The MADDPG actor into which its weights were copied received, from the simulator, the ego speed in metres per second, the SUMO lane index, an on-ramp flag and a count of nearby vehicles---different quantities, different scales and, for the lane, the opposite orientation (SUMO numbers lanes from the right, NGSIM from the left). The frozen reference used by the KL term was queried on the same raw tuple, and so was the safety evaluator; the NGSIM validator, by contrast, fed the same network the normalised features it had been cloned on. Tables~I, II and IV of the submitted paper therefore describe the network on one input space and Table~III on another, and in neither case the space the fine-tuning had seen in both roles. The repository recorded the problem itself: the warm-start routine's current docstring states that the legacy path ``places BC's first-layer weights on mismatched inputs'' and is kept for reproducibility, and the aligned projection was committed one and two days after the submission date. The only rows unaffected are those of the BC-only policy, which was cloned and scored in the same space. Every other method in Tables~I--V was scored on inputs its network had never been trained on, so none of those columns can be read as a statement about simulator behaviour.

\paragraph{Action semantics (D2).}
The prior emits probabilities in the NGSIM order (Lane\_Left, Lane\_Right, Accelerate, Decelerate, Maintain); the environment executes integers in the order (keep, accelerate, decelerate, left, right). No permutation existed before the fix committed after submission, whose message states that without it the KL term compares semantically mismatched distributions. During every submitted run the prior's Maintain mass was executed as Lane\_Right, Accelerate as Decelerate, Decelerate as Lane\_Left and Lane\_Right as Accelerate. Independently, the environment and the safety evaluator executed the simulator's ``left'' as a move to lane $k-1$, which in SUMO's right-to-left numbering is physically right. For a reinforcement learner alone the labels are immaterial; for a policy anchored to human demonstrations they are the content of the anchor. A KL penalty towards a fixed, high-entropy reference may still explain why the submitted CALM did not collapse, but that reference was not human-like in simulator semantics, and the claim that CALM anchors simulator behaviour to human driving is unsupported for the submitted runs.

\paragraph{Lane grid (D3).}
The projection $\phi_{\mathrm{BC}}$ that was added after submission to resolve D1 computed the lane feature as $(6-\ell)/6$ for SUMO lane $\ell$, which puts lane~0 at $1.0$. The dataset stores $(\mathrm{Lane\_ID}-1)/6$, so $1.0$ is Lane~7, the on-ramp, and the rightmost through lane, Lane~6, is $5/6$. Every simulator lane therefore queried the prior at the grid point of the lane one to its right. The submitted runs predate the projection and are unaffected; we list the defect because it would have carried into every corrected run had the contract not been checked feature by feature, and because it did carry into the anchor of our companion robustness study. The correction (Section~\ref{sec:contract_obs}) projects onto the dataset grid and makes the validator refuse any checkpoint whose recorded contract does not say so.

\subsection{Simulator evaluator (D4--D6)}
\label{sec:defects_evaluator}

\paragraph{Agents dropped before departure (D4).}
The safety evaluator initialised its tracking set with the five agent identifiers and removed any identifier absent from the simulation after the first step. Agents depart one second apart, so four of the five were removed before they appeared, and of the three exit-assigned agents only the first could ever be counted. This is the whole of the submitted Table~III's exit column: CALM's ``33.3\,\% exit success, exactly one exit per episode'' is the first agent exiting in every one of 100 episodes while the other two exit agents, also exiting in every episode, were not tracked. Re-measuring the same checkpoints under the same conditions with the lifecycle repaired (Section~\ref{sec:concealed}, Table~\ref{tab:t2_table3}) raised every method's exit success---CALM to 100\,\%, QMIX-CALM from 26.7 to 81.0\,\%, entropy-only from 20.0 to 60.0\,\%, DAgger from 7.7 to 28.3\,\%, GAIL from 6.7 to 20.0\,\%---while the speed and proximity columns, which were computed over every vehicle present rather than over the tracking set, reproduced the submitted values to full floating-point precision. The ordering of the methods did not change; every sentence that interpreted the magnitudes did.

\paragraph{SUMO drove the agents' lane changes (D5).}
Neither the training environment nor the evaluator set a lane-change mode for agent vehicles, so SUMO's own lane-change model, with its default mode, moved them autonomously whenever the policy did not. In the network only lane~0 of the first highway segment connects to the exit ramp; an exit agent that departs in lanes 1--3 reaches the ramp only if something moves it right. Under the default mode that something was frequently the simulator: the constant-Maintain policy, which never commands a lane change, exits in 100\,\% of episodes with the highest speed and the fewest near-misses of all pilot policies, and per-step logging showed that 70\,\% of the exit-lane moves of the unconstrained MADDPG policy and 48\,\% of those of the entropy-only policy were SUMO's. With autonomous agent lane changes disabled (mode 512, Section~\ref{sec:contract_lc}) the same constant policy exits in 7\,\% of episodes, those in which the first agent happened to depart in lane~0. Exit success under the submitted contract measured SUMO's driver model for every policy whose commanded share was low.

\paragraph{Collisions (D6).}
The evaluator defined a collision as a SUMO teleport of an agent vehicle. SUMO's car-following model and gap enforcement brake for every vehicle, including agents that command nothing or oscillate between lanes, and a teleport requires a prolonged blockage. The count was zero for every method in the submitted study, zero in the 500-episode re-measurement, zero in all 3{,}600 episodes we ran under the corrected contract, and zero by the direct-collision definition in the 2{,}900 of those in which it was counted---for stranded do-nothing agents and for policies changing lanes 500 times per episode alike. Zero collisions is a property of the simulator, and the submitted ``best safety profile'' rested on it together with D4 and D5.

\subsection{The NGSIM comparison (D7--D10)}
\label{sec:defects_ngsim}

\paragraph{In-sample evaluation (D7).}
The submitted paper described agreement on a ``held-out NGSIM test set''. The code contained no vehicle- or time-based split. Behavioural cloning used a random, action-stratified 90/10 split of the rows of the full 47{,}874-row, 567-vehicle file; the validator scored every row of that file for every method; and DAgger's interactive expert, a $k$-nearest-neighbour vote with $k=5$, was built on the same file it was later scored against. Consecutive rows of one vehicle are near-duplicates (nearest-neighbour feature distance within a vehicle is about six times smaller than between vehicles), so a row split does not hold anything out. The size of the effect for the most exposed number is measurable: the $k$-NN expert agrees with the human label in 77.2\,\% of rows when the evaluated row is in its tree, 64.7\,\% when only that row is removed, 43.2\,\% leave-one-vehicle-out and 41.7\,\% when built on the training vehicles and scored on the test vehicles. DAgger's 45.3\,\% was agreement with labels its own expert had memorised. The cloned prior does not over-fit (train and validation accuracy 35.3\,\% each), so its number barely moves under a true hold-out; the claim was still false for every method.

\paragraph{No majority baseline (D8).}
Maintain is 54.8\,\% of the labels in the full file and 55.3\,\% in the vehicle-grouped validation split; the two lane-change actions together are 0.26\,\%. Agreement is therefore a three-class problem with a trivial baseline, constant Maintain, that none of the submitted methods and none of the corrected ones exceeds (Section~\ref{sec:agreement}). The submitted paper reported no such row, and the ``human-like'' reading of agreement in the thirties and forties depended on its absence.

\paragraph{The entropy was not Definition~1 (D9).}
The paper defined mode collapse through the expected per-state entropy of $\pi(\cdot\,|\,o)$ (Definition~1) and proved a bound on that quantity. The reported ``action entropy'' was the entropy of the histogram of actions the policy selected across the dataset---with the selection made by sampling for GAIL and DAgger and by argmax for every other method. A deterministic policy whose argmax varies with the state scores high on this histogram while its per-state entropy is zero; a stochastic policy that ignores the state scores exactly its marginal. Recomputing all 62 saved checkpoints under one selection rule (argmax) and the three quantities of Section~\ref{sec:contract_entropy}, we found that CALM's 0.88 is 0.76 under Definition~1 (and the BC prior's 0.88 is 0.76, the two agreeing with each other as in the paper); that GAIL's and MASAC's selected-action entropy is 0.000 under argmax, their per-state entropies 0.50 and 0.56; that the entropy-only runs with $\alpha\le0.1$ pick one argmax for nearly every state yet have per-state entropies from 0.21 to 0.94, exceeding CALM's for $\alpha\ge0.1$; and that DAgger's agreement is 57.5\,\% under argmax against the 45.3\,\% reported from sampling. Whether those per-state values describe the trained policies or saturated outputs on inputs the networks never saw (D1) cannot be decided from these checkpoints, which is one more reason the collapse claims are re-established in Section~\ref{sec:results} rather than carried over. The statement ``Maintain 99.7\,\%, entropy 0.003'' for the unconstrained baseline is consistent under either definition; every other entropy in the submitted paper is a selected-histogram quantity.

\paragraph{One actor of five (D10).}
The validator scored \texttt{actors[0]} of each five-agent checkpoint. The agreement and entropy columns of Table~I therefore describe one actor; the corrected validator scores every actor, reports the mean and the spread across actors, and computes entropies per actor before averaging, because pooling deterministic actors' choices would read their disagreement as the diversity of one policy.

\subsection{Reporting (D11)}
\label{sec:defects_reporting}

``Final reward'' was the mean over the last ten training episodes, with exploration noise, and the evaluator records no reward, so the paper's reward comparison is between training statistics, not evaluation outcomes; the corrected study keeps the quantity, labelled as such (Section~\ref{sec:reward}). Several protocol statements were also inexact: not every method was run on the stated five seeds, QMIX-CALM was run with ten agents while the paper stated five, and the imitation trainer saved its last-epoch rather than its best-validation weights (the snapshot aliased the live tensors), carried no seed, and computed its class-weighted validation loss in a batch-dependent way that could change the early-stopping epoch. None of these alters the conclusions of this section; all are corrected in the pipeline of Section~\ref{sec:contract} and in the historical record of the appendix.

\subsection{What the defects did together}
\label{sec:defects_together}

Each defect pushed in the direction of the submitted conclusions. D1 and D2 broke the link between the anchor and human behaviour while leaving a high-entropy reference in place, so the method did not collapse but did not anchor to anything human either. D4 understated exits for every method, which turned CALM's ``one exit per episode'' into a result worth explaining and hid that two constant policies exit in every episode. D5 and D6 gave every policy the simulator's lane changes and the simulator's collision avoidance, so the safety table ranked policies by how little they interfered with SUMO. D7 removed the hold-out, D8 the baseline that would have placed every agreement figure below chance-with-priors, and D9 reported a quantity on which a two-action greedy policy looks diverse. Under the corrected contract, as Section~\ref{sec:results} shows, the greedy anchored policy is the constant Lane\_Right policy, and the submitted paper's safety and human-likeness claims are that policy's and the simulator's.

% T1: hand-written (not produced by make_r2_tables.py). Sources: VANET protocol §1, §4b–§4d, §8.7–§8.8; ESWA plan §2b–§2d, §9.
\begin{table*}[!t]
\centering
\footnotesize
\caption{Defects of the submitted pipeline, their mechanism, their measured effect on the reported numbers and their correction. D1--D3 are defects of the input/action contract, D4--D6 of the simulator evaluator, D7--D10 of the NGSIM comparison, D11 of the reporting. Section numbers refer to this paper.}
\label{tab:t1_defects}
\setlength{\tabcolsep}{4pt}
\begin{tabular}{@{}>{\raggedright\arraybackslash}p{0.03\textwidth}>{\raggedright\arraybackslash}p{0.14\textwidth}>{\raggedright\arraybackslash}p{0.31\textwidth}>{\raggedright\arraybackslash}p{0.27\textwidth}>{\raggedright\arraybackslash}p{0.185\textwidth}@{}}
\toprule
\# & Defect & Mechanism & Measured effect & Correction (Section~\ref{sec:contract}) \\
\midrule
D1 & Prior and actor consumed different feature spaces & The BC prior was fitted to four NGSIM-normalised features in $[0,1]$; the actor it was copied into received raw simulator quantities (speed in m/s, SUMO lane index, on-ramp flag, neighbour count), as did the frozen KL reference and the safety evaluator; the NGSIM validator fed the same network normalised features. & Every number in Tables~I--V of the submitted paper describes one network scored on an input space it was not trained on, except the BC-only rows. The KL anchor pulled towards a reference evaluated off-distribution. & One observation schema for training, anchor and both evaluators (the $\phi_{\mathrm{BC}}$ projection, \S\ref{sec:contract_obs}). \\
D2 & Action semantics unpermuted and lateral direction inverted & The BC/NGSIM order (Lane\_Left, Lane\_Right, Accelerate, Decelerate, Maintain) was never mapped to the simulator order (keep, accel, decel, left, right): the prior's Maintain executed as Lane\_Right, Accelerate as Decelerate, Decelerate as Lane\_Left. The environment also executed ``left'' as SUMO lane $k-1$, which is physically right. & The anchor was not human-like in simulator semantics; ``anchors behaviour to human demonstrations'' is unsupported for the submitted runs. & One action order with an explicit index map at the environment boundary; left executes toward the median; tests pin both (\S\ref{sec:contract_obs}). \\
D3 & Lane feature one NGSIM lane off & The projection added after submission to repair D1 computed $(6-\ell)/6$, placing SUMO lane 0 at the dataset value of the on-ramp (Lane~7); the dataset stores $(\mathrm{Lane\_ID}-1)/6$. Every simulator lane queried the prior at the grid point of the lane to its right. & No submitted number (the submitted runs had no projection); it would have carried into every corrected run and did carry into the companion study's anchor. & Projection onto the dataset grid; the validator refuses a checkpoint whose contract does not declare it (\S\ref{sec:contract_obs}). \\
D4 & Agents dropped before departure & The safety evaluator discarded any agent absent after the first simulation step; agents depart at $t=0,\dots,4$\,s, so four of five were never tracked and only one exit agent could be counted. & Exit success under-counted for every method: CALM 33.3\,\% $\to$ 100\,\%, QMIX-CALM 26.7 $\to$ 81.0, entropy-only 20.0 $\to$ 60.0, DAgger 7.7 $\to$ 28.3, GAIL 6.7 $\to$ 20.0 (Table~\ref{tab:t2_table3}); speed and proximity unchanged to full precision. & Pending $\to$ active $\to$ exited lifecycle; departed agents counted (\S\ref{sec:contract_safety}). \\
D5 & SUMO drove the agents' lane changes & No lane-change mode was set for agent vehicles, so SUMO's lane-change model (mode 1621) moved them autonomously; a policy that never commands a lane change still exits. & Constant Maintain reaches 100\,\% exits, the highest speed and the fewest near-misses of all pilot policies; 70\,\% of unconstrained MADDPG's and 48\,\% of entropy-only's exit-lane moves were SUMO's. Under mode 512 the same constant policy exits 7\,\%. & SUMO lane-change mode 512 for agent vehicles; commanded and uncommanded transitions counted apart (\S\ref{sec:contract_lc}). \\
D6 & Collision defined as teleport; zero read as a policy property & SUMO's car-following and gap enforcement brake for every agent; a teleport requires a prolonged blockage. & Zero teleports in all $\approx$\,4{,}100 evaluation episodes under both contracts and zero direct collisions in the 3{,}200 in which they were counted, oscillating and stranded policies included. ``Best safety profile'' rested on a simulator property. & Direct collisions and teleports reported apart and interpreted as the simulator's; exposure-normalised proximity beside reference policies (\S\ref{sec:contract_safety}). \\
D7 & In-sample NGSIM comparison & BC used a random 90/10 row split of the whole 47{,}874-row file and the validator scored the whole file; DAgger's $k$-NN expert was built on the whole file. & The $k$-NN expert agrees with the labels at 77.2\,\% in-sample and 43.2\,\% leave-one-vehicle-out; ``held-out NGSIM test set'' was false for every method. & Vehicle-grouped 70/15/15 split with leakage checks; all selection on \texttt{val.csv}; \texttt{test.csv} frozen (\S\ref{sec:contract_data}). \\
D8 & No majority-class baseline & Maintain is 55.3\,\% of the labels and lane changes 0.26\,\%; agreement is a three-class problem with a trivial baseline. & Every method, submitted or corrected, scores below constant Maintain (Table~\ref{tab:t4_val}). & Majority row in every agreement table (\S\ref{sec:contract_refs}). \\
D9 & ``Action entropy'' was not Definition~1 & The reported quantity was the entropy of the histogram of selected actions over the dataset, with sampling for GAIL and DAgger and argmax for the rest. & CALM's 0.88 is 0.76 under Definition~1; GAIL's and MASAC's selected-action entropy is 0.000 under one rule; DAgger's 45.3\,\% agreement was sampled, 57.5\,\% under argmax. & Three named entropy quantities from the softmax output and one declared selection rule (\S\ref{sec:contract_entropy}). \\
D10 & One actor of five scored & The validator scored \texttt{actors[0]} of each five-agent checkpoint. & Table~I's agreement and entropy describe one actor. & Every actor scored; mean and spread across actors reported; entropies computed per actor before averaging (\S\ref{sec:contract_entropy}). \\
D11 & Training-time reward reported as performance; protocol statements inexact & ``Final reward'' was the mean of the last ten training episodes with exploration noise; the evaluator records no reward. Seeds, run counts and QMIX agent count were not as run. & Reward columns compare training statistics; the stated five seeds and $n=5$ did not hold throughout. & Reward labelled as a training statistic; seeds, counts and contract stated per table (\S\ref{sec:contract_report}). \\
\bottomrule
\end{tabular}
\end{table*}


% Section 4: corrected evaluation contract.
% Evidence: VANET protocol §2–§5, §8.1–§8.5, §8.8 preamble, §8.9 (design, evaluation, bootstrap), §8.11 preamble.
\section{Corrected Evaluation Contract}
\label{sec:contract}

Every number in Section~\ref{sec:results} was produced under the contract below. The contract is recorded in each training run's manifest, in the header of its training log and in every checkpoint it writes (observation schema, action order, lane direction, lane grid, lane-change mode, action-selection rule, git commit and seed), and both evaluators refuse a checkpoint whose recorded contract does not match the one they are asked to apply. The unit tests that pin each item run without the simulator.

\subsection{One observation space and one action order}
\label{sec:contract_obs}

\paragraph{Observation.}
The policy input in the simulator \emph{is} the projection $\phi_{\mathrm{BC}}$ onto the four features the prior was cloned on: lane position, velocity, distance to the exit and headway, each normalised to $[0,1]$ as in the dataset. The lane feature uses the dataset's grid: SUMO lane $k$ of the four-lane segment is mapped to NGSIM Lane~$6-k$ and encoded as $(\mathrm{Lane\_ID}-1)/6$, so the four simulator lanes read as Lanes 6 to 3, the rightmost through lane of the recording and the three to its left. The cloned prior, the warm-started actor, the frozen KL reference, the NGSIM validator and the safety evaluator all consume this one vector. The reward is computed from the raw simulator state, not from the policy input.

\paragraph{Action order.}
Actors emit probabilities in the NGSIM order; the environment converts each sampled or argmax index to the simulator's action with a fixed map at its boundary, so the KL term compares like with like without any permutation, and the NGSIM validator compares the actor's vector with the human label directly. ``Left'' executes toward the median. Tests pin that the index maps are mutual inverses and that the prior's Maintain executes as SUMO's keep.

\paragraph{Warm start.}
When a prior is given, its weights are copied into every actor and every target network before the first episode, so an anchored and an unanchored run can start from identical parameters.

\subsection{Data: a vehicle-grouped split with a frozen test set}
\label{sec:contract_data}

The 47{,}874 rows (567 vehicles) of the discretised Interstate~80 file were split by vehicle into training, validation and test sets of 70/15/15\,\% of the vehicles, stratified by each vehicle's dominant action and row-count tercile: 396 vehicles (31{,}476 rows), 87 (8{,}143) and 84 (8{,}255). An automated check verifies zero vehicle overlap and zero identical feature rows between any two splits, records the SHA-256 of each file in a manifest against which every consumer is checked, and reports action-share parity (maximum deviation from the full file 1.7 points, on the test split) and nearest-neighbour feature distances (median within-train leave-one-out 0.0022; validation-to-train 0.0128; test-to-train 0.0147). Because the file comes from one recording, 91\,\% of validation and 77\,\% of test vehicles overlap some training vehicle's time window; features are ego-centric and we state this rather than claim temporal independence. The validation split has Maintain 55.3\,\%, Decelerate 22.9\,\%, Accelerate 21.6\,\%, Lane\_Right 0.22\,\% and Lane\_Left 0.01\,\%.

The prior is cloned on the training split with early stopping on the validation split; every learning curve, every checkpoint score and every decision in this paper uses the validation split; the test split is frozen. Its labels were read once, before the evaluation rules were fixed, for the $k$-NN leave-one-vehicle-out diagnostic reported in Section~\ref{sec:defects_ngsim}; that access is logged and no model, hyper-parameter or protocol choice depended on it. The claim available to this paper is therefore ``no model selection consulted the test split'', not ``untouched''. \todo{If the one-time test-split scoring of the final checkpoints is approved, state here that it was performed once, after every table was final, and name the tables; otherwise delete this note and leave the agreement tables on \texttt{val.csv}.}

\subsection{The prior}
\label{sec:contract_prior}

The prior is a \emph{class-balanced} imitation policy, as in the submitted method: the cross-entropy loss uses inverse-frequency class weights, so lane changes, 0.26\,\% of the labels, are up-weighted by two to three orders of magnitude relative to Maintain. Its probabilities are therefore not human action frequencies. On the training split its argmax is Lane\_Right in 18.8\,\% of states against 0.2\,\% for the humans, and Maintain in 26.3\,\% against 54.3\,\%; on the validation split it agrees with the human label in 34.85\,\% of states, has Jensen--Shannon divergence 0.269 from the human marginal and per-state entropy 0.775. The trainer saves the weights of the epoch with the lowest class-weighted validation loss (computed batch-invariantly as $\sum w\ell / \sum w$), under a fixed seed, and records the split hashes, the seed, the weight vector and the restored epoch in the checkpoint. An \emph{unweighted} prior, trained identically without class weights, is a declared ablation: it matches the human marginal (Jensen--Shannon 0.012) and its argmax is Maintain in 94\,\% of states.

\subsection{Three entropy quantities and one selection rule}
\label{sec:contract_entropy}

For a policy with softmax output $\pi(\cdot\,|\,o)$ over the five actions and a set of observations $\mathcal{O}$ we report, each normalised by $\log 5$:
\begin{itemize}
  \item the \emph{per-state entropy} $\frac{1}{|\mathcal{O}|}\sum_{o\in\mathcal{O}} \mathcal{H}(\pi(\cdot\,|\,o))$, the estimator of Definition~1;
  \item the entropy of the \emph{mean probability vector} $\mathcal{H}\big(\frac{1}{|\mathcal{O}|}\sum_o \pi(\cdot\,|\,o)\big)$;
  \item the \emph{selected-action entropy}, the entropy of the histogram of actions chosen under the declared selection rule---the quantity the submitted paper reported.
\end{itemize}
Probabilities are the actors' softmax outputs; no second softmax is applied. On the NGSIM validation split $\mathcal{O}$ is the set of human states and the selection rule is argmax, for every policy, the prior included; over simulator rollouts $\mathcal{O}$ is the set of agent-steps, the policy's probability vector being logged at each step, so Definition~1 is estimated on the states the policy itself visits under the declared execution rule. Every one of the five actors is scored; agreement, divergence and the three entropies are reported as the mean across actors with the spread beside it, each entropy computed per actor before averaging. For QMIX, which executes greedily, the ``per-state entropy'' is that of $\mathrm{softmax}(Q/\tau)$, a proxy for the spread of action values and not the entropy of an execution policy; the selected-action entropy is the behaviourally meaningful one for QMIX and is labelled so.

\subsection{A declared execution rule}
\label{sec:contract_exec}

MADDPG's action selection during training is \emph{sampled} from the actor's softmax whenever a BC coefficient, an entropy coefficient or a non-default temperature is set, and greedy with $\epsilon$-exploration otherwise---so in the corrected study's first runs (Section~\ref{sec:collapse}) the anchored and the entropy-regularised methods sampled while unconstrained MADDPG did not, and the submitted paper never stated the rule. The corrected code decouples the rule from the regularisers (\texttt{--act-sampling} sampled or greedy) and records it. At evaluation the rule is declared separately (\texttt{--execution} sampled or greedy); \emph{sampled} is the rule the anchored policy trains with and is the primary rule in Section~\ref{sec:results}, and greedy results are reported apart and never mixed with sampled ones in a table. Section~\ref{sec:concealed} shows how far the two rules diverge for the same checkpoint.

\subsection{Policy-attributable lane changes}
\label{sec:contract_lc}

Agent vehicles are inserted with SUMO lane-change mode 512: no autonomous lane change, while a lane change requested by the policy through TraCI still respects other drivers' gaps, so the simulator's safety layer is kept and only the decision to move is the policy's. Of the first segment's four lanes only lane~0 connects to the exit ramp, so under this mode an exit agent that departs in lanes 1--3 reaches the ramp only if the policy commands the move; exit success becomes a statement about the policy. The evaluator logs every lane transition and whether it was commanded, and reports commanded and uncommanded transitions, left and right, together with the share of mainline time spent in lane~0. Both the training environment and the evaluator apply the mode, and the mode is part of the recorded contract.

\subsection{Safety accounting}
\label{sec:contract_safety}

The evaluator tracks each agent from pending through active to exited and reports how many departed. \emph{Exit success} is the fraction of the three exit-assigned agents whose last edge was a ramp edge when they left the network; exits per episode (0--3) is reported beside it. \emph{Direct collisions} (SUMO's collision list) and \emph{teleports} are counted apart, and both are interpreted as properties of the simulator's car-following layer rather than of the policy (Section~\ref{sec:concealed}). A \emph{proximity event} is an agent--vehicle pair closer than 5\,m at a simulation step, summed over agents and steps; it is reported as a raw count per episode, per 1{,}000 agent-steps and per agent-kilometre, the last as the pooled rate $\sum\text{events}/\sum\text{agent-km}$ because slow and stranded policies accumulate steps without distance. Average speed is computed over every agent present in the simulation. Every simulator figure in this paper states the lane-change mode and the execution rule under which it was measured.

\subsection{Reference policies}
\label{sec:contract_refs}

Three policies that learn nothing are run through the same evaluator on the same episodes: \emph{constant Maintain}, \emph{constant Lane\_Right} and \emph{SUMO default}, in which no action is ever commanded. Under the default lane-change mode the first and the third coincide; under mode 512 constant Maintain and SUMO default are do-nothing policies and constant Lane\_Right is the optimal exit policy of this network. On the NGSIM side the \emph{majority class}, constant Maintain, is a row in every agreement table (55.3\,\% on the validation split), with the uniform-random and the submitted rule-based policy beside it.

\subsection{Matched episodes, checkpoints and inference}
\label{sec:contract_stats}

Each training run is evaluated at its final checkpoint (episode 300 for MADDPG, 100 for QMIX); checkpoints every 50 episodes are scored on the validation split for learning curves only, and no checkpoint is selected on any score. Simulator evaluation uses 50 episodes per checkpoint with episode seed $s\times10^4+e$ for training seed $s$ and episode index $e$, so every arm, every reference policy and the frozen prior see the same 50 traffic realisations for a given training seed. Where two policies are compared, the comparison is paired over episodes within a training seed: 10{,}000 bootstrap resamples of matched episode pairs, with the proximity rate recomputed as a pooled ratio inside each resample, give 95\,\% intervals for differences in entropy, exits and speed and for the ratio of proximity rates. An episode is the unit; vehicle-steps are not treated as independent. The intervals describe the evaluation uncertainty of the trained checkpoints and do not remove the limitation that the isolation of Section~\ref{sec:anchor} rests on two training seeds. The decision rule for that isolation---thresholds, guards and the wording each outcome permits---was written and reviewed before any of its runs and is reproduced in Section~\ref{sec:anchor}.

\subsection{Reporting rules}
\label{sec:contract_report}

Reward is a training statistic: the mean over the last 50 training episodes, with exploration, under the arm's training rule; the evaluator returns none, and no reward figure is presented as an evaluation outcome. Seeds, episode counts, agent counts and the contract items above are stated for every table. Nothing measured on the submitted pipeline is used as a comparator for the corrected runs: the historical numbers appear only in Section~\ref{sec:concealed}, beside their re-measurement, and in the appendix, labelled as measurements of the defective pipeline.


\section{Results}
\label{sec:results}

\subsection{What the defects concealed}
\label{sec:concealed}
\todo{Protocol \S4c and \S8.8 around Tables~\ref{tab:t2_table3} and \ref{tab:t3_gate}: the exit-counting artefact; baseline identities; execution-rule dependence; zero collisions in 3,600 episodes.}
\begin{table*}[t]
\centering
\caption{Submitted Table III (historical: defective pipeline, legacy observation contract) beside its re-measurement with the agents tracked for the whole episode (protocol \S4c; same checkpoints, same legacy contract, 100 episodes per method). The submitted exit column counted only agents still tracked after step 1.}
\label{tab:t2_table3}
\footnotesize
\resizebox{\textwidth}{!}{%
\begin{tabular}{lccccccc}
\toprule
\textbf{Method} & \textbf{Exit \% (subm.)} & \textbf{Exit \% (re-meas.)} & \textbf{Speed m/s (subm.)} & \textbf{Speed m/s (re-meas.)} & \textbf{Prox./ep (subm.)} & \textbf{Prox./ep (re-meas.)} & \textbf{Collisions \%} \\
\midrule
CALM (class-balanced prior) & 33.3 & 100.0 & 25.8 $\pm$ 1.2 & 25.8 $\pm$ 1.2 & 42.5 & 42.5 & 0.0 \\
QMIX-CALM ($n=5$) & 26.7 & 81.0 & 19.7 $\pm$ 6.9 & 19.7 $\pm$ 6.9 & 220.9 & 220.9 & 0.0 \\
Entropy-only ($\alpha=0.01$) & 20.0 & 60.0 & 16.3 $\pm$ 6.1 & 16.3 $\pm$ 6.1 & 685.4 & 685.4 & 0.0 \\
DAgger & 7.7 & 28.3 & 11.0 $\pm$ 3.6 & 11.0 $\pm$ 3.6 & 3150.5 & 3150.5 & 0.0 \\
GAIL & 6.7 & 20.0 & 8.1 $\pm$ 2.2 & 8.1 $\pm$ 2.2 & 2563.8 & 2563.8 & 0.0 \\
\bottomrule
\end{tabular}}
\vspace{1mm}
\parbox{\textwidth}{\scriptsize Historical record only: these checkpoints were trained on raw SUMO features with an unpermuted action order (protocol \S1) and are not comparators for any corrected-contract result. Collision rate as submitted = fraction of episodes with $\geq 1$ teleport; 0 in every cell.}
\end{table*}

\begin{table*}[t]
\centering
\caption{Diagnostic gate (protocol \S8.8): the same checkpoints evaluated under SUMO's default lane changing and under policy-attributable lane changes (\texttt{agent\_lane\_change\_mode=512}), with greedy and sampled execution, beside the three reference policies. 50 episodes per seed; zero direct collisions and zero teleports in every cell.}
\label{tab:t3_gate}
\footnotesize
\resizebox{\textwidth}{!}{%
\begin{tabular}{lcccccccc}
\toprule
\textbf{Contract, execution: policy} & \textbf{$n$} & \textbf{Exit \%} & \textbf{Speed m/s} & \textbf{Prox./agent-km} & \textbf{$H$ per-state} & \textbf{$H$ selected} & \textbf{Lane changes L / R} & \textbf{Lane-0 share} \\
\midrule
SUMO default, greedy: CALM (class-balanced prior) & 100 & 100.0 & 25.3 & 7.5 & 0.65 & 0.01 & 555 / 0 & 0.96 \\
SUMO default, greedy: CALM (unweighted prior) & 100 & 100.0 & 27.2 & 6.2 & 0.44 & 0.00 & 0 / 553 & 0.50 \\
SUMO default, greedy: Entropy-only ($\alpha=0.01$) & 100 & 82.3 & 18.9 & 56.2 & 0.36 & 0.74 & 280 / 260 & 0.62 \\
SUMO default, greedy: MADDPG, no prior & 50 & 100.0 & 26.8 & 11.7 & 0.36 & 0.81 & 84 / 193 & 0.52 \\
SUMO default, greedy: QMIX-CALM ($n=5$) & 50 & 100.0 & 24.4 & 10.6 & -- & 0.33 & 278 / 0 & 0.96 \\
SUMO default, greedy: Constant Maintain & 100 & 100.0 & 27.2 & 6.2 & 0.00 & 0.00 & 0 / 553 & 0.50 \\
SUMO default, greedy: Constant Lane\_Right & 100 & 100.0 & 25.3 & 7.5 & 0.00 & 0.00 & 555 / 0 & 0.96 \\
SUMO default, greedy: SUMO default (no commands) & 100 & 100.0 & 27.2 & 6.2 & -- & 0.00 & 0 / 553 & 0.50 \\
\midrule
SUMO default, sampled: CALM (class-balanced prior) & 100 & 79.7 & 12.1 & 37.2 & 0.71 & 0.72 & 553 / 12 & 0.93 \\
SUMO default, sampled: CALM (unweighted prior) & 100 & 97.3 & 16.7 & 21.2 & 0.43 & 0.43 & 440 / 332 & 0.69 \\
SUMO default, sampled: Entropy-only ($\alpha=0.01$) & 100 & 80.3 & 18.3 & 40.4 & 0.31 & 0.88 & 565 / 436 & 0.46 \\
SUMO default, sampled: MADDPG, no prior & 50 & 100.0 & 25.1 & 17.2 & 0.36 & 0.85 & 395 / 293 & 0.48 \\
\midrule
mode 512, greedy: CALM (class-balanced prior) & 100 & 100.0 & 24.5 & 8.3 & 0.65 & 0.01 & 555 / 0 & 0.96 \\
mode 512, greedy: CALM (unweighted prior) & 100 & 7.0 & 8.4 & 681.0 & 0.41 & 0.00 & 0 / 0 & 0.01 \\
mode 512, greedy: Entropy-only ($\alpha=0.01$) & 100 & 49.7 & 11.8 & 62.4 & 0.31 & 0.84 & 277 / 0 & 0.25 \\
mode 512, greedy: MADDPG, no prior & 50 & 33.3 & 12.2 & 254.0 & 0.44 & 0.57 & 84 / 0 & 0.09 \\
mode 512, greedy: QMIX-CALM ($n=5$) & 50 & 100.0 & 23.4 & 11.6 & -- & 0.33 & 278 / 0 & 0.96 \\
mode 512, greedy: Constant Maintain & 100 & 7.0 & 8.4 & 681.0 & 0.00 & 0.00 & 0 / 0 & 0.01 \\
mode 512, greedy: Constant Lane\_Right & 100 & 100.0 & 24.5 & 8.3 & 0.00 & 0.00 & 555 / 0 & 0.96 \\
mode 512, greedy: SUMO default (no commands) & 100 & 7.0 & 8.4 & 681.0 & -- & 0.00 & 0 / 0 & 0.01 \\
\midrule
mode 512, sampled: CALM (class-balanced prior) & 100 & 78.0 & 11.1 & 41.4 & 0.71 & 0.72 & 563 / 0 & 0.94 \\
mode 512, sampled: CALM (unweighted prior) & 100 & 95.0 & 15.0 & 33.8 & 0.43 & 0.43 & 769 / 2 & 0.53 \\
mode 512, sampled: Entropy-only ($\alpha=0.01$) & 100 & 65.0 & 15.9 & 45.0 & 0.26 & 0.91 & 845 / 1 & 0.23 \\
mode 512, sampled: MADDPG, no prior & 50 & 100.0 & 23.8 & 29.7 & 0.36 & 0.84 & 599 / 2 & 0.31 \\
\bottomrule
\end{tabular}}
\vspace{1mm}
\parbox{\textwidth}{\scriptsize Per-state entropy is the mean normalised entropy of the policy distribution over visited states; selected-action entropy is the normalised entropy of the executed-action histogram. Under SUMO default, greedy class-balanced-prior CALM reproduces constant Lane\_Right and greedy unweighted-prior CALM reproduces constant Maintain and the SUMO default to the decimal. Under mode 512 the stranded reference policies (constant Maintain, SUMO default) exit 7\,\% and expose two orders of magnitude more proximity per km. QMIX-CALM has no per-state entropy (argmax policy).}
\end{table*}


\subsection{Agreement against the majority baseline}
\label{sec:agreement}
\todo{Protocol \S8.1a, \S8.6, \S8.7, \S8.10 around Table~\ref{tab:t4_val}.}
\begin{table*}[t]
\centering
\caption{Action agreement with the held-out NGSIM validation split (\texttt{val.csv}, vehicle-grouped; protocol \S4) and Definition~1 entropies, all five actors scored, final checkpoints. The majority-class row is the trivial baseline no method reaches. \texttt{test.csv} has not been read.}
\label{tab:t4_val}
\footnotesize
\resizebox{\textwidth}{!}{%
\begin{tabular}{lcccccc}
\toprule
\textbf{Policy} & \textbf{Agreement \%} & \textbf{$\pm$ actor} & \textbf{JS} & \textbf{$H$ per-state} & \textbf{$H$ selected} & \textbf{Argmax share LL/LR/Acc/Dec/Maint} \\
\midrule
Majority class (constant Maintain) & 55.3 & -- & -- & 0.000 & 0.000 & 0.00/0.00/0.00/0.00/1.00 \\
Random & 20.0 & -- & 0.420 & 1.000 & 1.000 & 0.19/0.21/0.20/0.20/0.20 \\
Rule-based controller & 22.4 & -- & 0.486 & 0.000 & 0.641 & 0.00/0.29/0.00/0.50/0.21 \\
\midrule
BC prior, class-balanced (Stage 0) & 34.9 & -- & 0.269 & 0.775 & 0.878 & 0.01/0.21/0.25/0.21/0.32 \\
BC prior, unweighted (\S8.7) & 56.2 & -- & 0.012 & 0.615 & 0.133 & 0.00/0.00/0.00/0.06/0.94 \\
\midrule
CALM pilot, seed 42 & 38.7 & 3.0 & 0.302 & 0.808 & 0.374 & 0.00/0.30/0.00/0.00/0.70 \\
CALM pilot, seed 123 & 39.6 & 2.0 & 0.299 & 0.806 & 0.365 & 0.00/0.28/0.00/0.00/0.72 \\
CALM unweighted prior, seed 42 & 55.3 & 0.0 & 0.219 & 0.372 & 0.000 & 0.00/0.00/0.00/0.00/1.00 \\
CALM unweighted prior, seed 123 & 55.3 & 0.0 & 0.216 & 0.378 & 0.000 & 0.00/0.00/0.00/0.00/1.00 \\
Entropy-only, seed 42 & 4.6 & 8.5 & 0.622 & 0.506 & 0.089 & 0.14/0.66/0.21/0.00/0.00 \\
Entropy-only, seed 123 & 26.8 & 24.7 & 0.505 & 0.469 & 0.000 & 0.00/0.40/0.00/0.20/0.40 \\
MADDPG no prior, seed 42 & 9.1 & 10.3 & 0.584 & 0.563 & 0.020 & 0.20/0.39/0.40/0.00/0.01 \\
QMIX-CALM $n=5$, seed 42 (100 ep.) & 50.9 & 2.6 & 0.243 & 0.806 & 0.271 & 0.00/0.09/0.00/0.05/0.86 \\
\midrule
Option B B1 anchored, sampled-trained, seed 42 & 38.6 & 1.3 & 0.304 & 0.809 & 0.378 & 0.00/0.30/0.00/0.00/0.70 \\
Option B B1 anchored, sampled-trained, seed 123 & 39.8 & 3.5 & 0.300 & 0.806 & 0.362 & 0.00/0.28/0.00/0.00/0.72 \\
Option B B2 anchored, greedy-trained, seed 42 & 39.7 & 2.3 & 0.296 & 0.807 & 0.364 & 0.00/0.28/0.00/0.00/0.72 \\
Option B B2 anchored, greedy-trained, seed 123 & 39.3 & 2.1 & 0.297 & 0.806 & 0.369 & 0.00/0.28/0.00/0.00/0.72 \\
Option B B3 unanchored, sampled-trained, seed 42 & 16.2 & 5.7 & 0.440 & 0.632 & 0.297 & 0.07/0.29/0.44/0.14/0.05 \\
Option B B3 unanchored, sampled-trained, seed 123 & 33.9 & 16.7 & 0.377 & 0.450 & 0.175 & 0.00/0.10/0.29/0.21/0.40 \\
Option B B4 unanchored, greedy-trained, seed 42 & 16.4 & 5.7 & 0.436 & 0.623 & 0.325 & 0.06/0.30/0.46/0.13/0.05 \\
Option B B4 unanchored, greedy-trained, seed 123 & 26.3 & 12.9 & 0.384 & 0.579 & 0.283 & 0.04/0.15/0.39/0.21/0.21 \\
\bottomrule
\end{tabular}}
\vspace{1mm}
\parbox{\textwidth}{\scriptsize Agreement = share of validation rows on which the policy's argmax equals the human action. $\pm$ actor = standard deviation across the five actors. JS = Jensen--Shannon divergence between the policy's argmax histogram and the human action histogram. Lane changes are 0.26\,\% of human rows, so agreement is effectively a three-class problem.}
\end{table*}


\subsection{Collapse under the corrected contract}
\label{sec:collapse}
\todo{Pilot no-BC arms end as constant-action actors, a different constant per agent (\S8.6); the unweighted prior anchors CALM to constant Maintain (\S8.7): ``preserves the prior'' cuts both ways.}

\subsection{Isolating the anchor (pre-registered)}
\label{sec:anchor}
\todo{Protocol \S8.9--8.10 around Tables~\ref{tab:t5_option_b} and \ref{tab:t5b_decision}; Fig.~F2 entropy trajectories on val.csv; Fig.~F3 lane transitions and lane-0 occupancy.}
\begin{table*}[t]
\centering
\caption{Option B (protocol \S8.9--8.10): the same class-balanced prior fine-tuned for 300 episodes with and without the KL anchor, under sampled or greedy training action selection, two seeds, evaluated under \texttt{agent\_lane\_change\_mode=512} with both execution rules on matched episode seeds (50 per seed). Zero direct collisions and zero teleports in all 800 episodes.}
\label{tab:t5_option_b}
\footnotesize
\resizebox{\textwidth}{!}{%
\begin{tabular}{lcccccccc}
\toprule
\textbf{Evaluation execution: arm} & \textbf{$n$} & \textbf{Exit \%} & \textbf{Speed m/s} & \textbf{Prox./agent-km} & \textbf{$H$ per-state} & \textbf{$H$ selected} & \textbf{Exit agents' R moves: policy / SUMO} & \textbf{Lane-0 share} \\
\midrule
sampled: B1 anchored, sampled-trained, seed 42 & 50 & 85.3 & 11.9 & 37.4 & 0.714 & 0.723 & 280 / 0 & 0.95 \\
sampled: B1 anchored, sampled-trained, seed 123 & 50 & 80.0 & 11.2 & 41.0 & 0.706 & 0.717 & 278 / 0 & 0.95 \\
sampled: B1 anchored, sampled-trained, pooled & 100 & 82.7 & 11.6 & 39.2 & 0.710 & 0.720 & 558 / 0 & 0.95 \\
sampled: B2 anchored, greedy-trained, seed 42 & 50 & 85.3 & 11.8 & 38.3 & 0.720 & 0.730 & 282 / 0 & 0.96 \\
sampled: B2 anchored, greedy-trained, seed 123 & 50 & 69.3 & 10.4 & 44.5 & 0.707 & 0.718 & 278 / 0 & 0.95 \\
sampled: B2 anchored, greedy-trained, pooled & 100 & 77.3 & 11.1 & 41.4 & 0.713 & 0.724 & 560 / 0 & 0.95 \\
sampled: B3 unanchored, sampled-trained, seed 42 & 50 & 38.7 & 12.1 & 482.9 & 0.254 & 0.774 & 3299 / 1 & 0.21 \\
sampled: B3 unanchored, sampled-trained, seed 123 & 50 & 45.3 & 10.5 & 121.9 & 0.216 & 0.758 & 258 / 0 & 0.35 \\
sampled: B3 unanchored, sampled-trained, pooled & 100 & 42.0 & 11.3 & 316.7 & 0.235 & 0.766 & 3557 / 1 & 0.28 \\
sampled: B4 unanchored, greedy-trained, seed 42 & 50 & 41.3 & 12.6 & 425.2 & 0.236 & 0.766 & 4447 / 0 & 0.23 \\
sampled: B4 unanchored, greedy-trained, seed 123 & 50 & 35.3 & 10.6 & 102.8 & 0.310 & 0.928 & 13071 / 0 & 0.26 \\
sampled: B4 unanchored, greedy-trained, pooled & 100 & 38.3 & 11.6 & 276.2 & 0.273 & 0.847 & 17518 / 0 & 0.24 \\
\midrule
greedy: B1 anchored, sampled-trained, seed 42 & 50 & 100.0 & 24.5 & 8.3 & 0.650 & 0.010 & 278 / 0 & 0.96 \\
greedy: B1 anchored, sampled-trained, seed 123 & 50 & 100.0 & 24.5 & 8.2 & 0.636 & 0.014 & 277 / 0 & 0.97 \\
greedy: B1 anchored, sampled-trained, pooled & 100 & 100.0 & 24.5 & 8.3 & 0.643 & 0.012 & 555 / 0 & 0.97 \\
greedy: B2 anchored, greedy-trained, seed 42 & 50 & 100.0 & 24.5 & 8.3 & 0.666 & 0.020 & 278 / 0 & 0.96 \\
greedy: B2 anchored, greedy-trained, seed 123 & 50 & 100.0 & 24.5 & 8.2 & 0.658 & 0.013 & 277 / 0 & 0.97 \\
greedy: B2 anchored, greedy-trained, pooled & 100 & 100.0 & 24.5 & 8.3 & 0.662 & 0.017 & 555 / 0 & 0.97 \\
greedy: B3 unanchored, sampled-trained, seed 42 & 50 & 6.0 & 9.8 & 554.6 & 0.389 & 0.730 & 117 / 0 & 0.17 \\
greedy: B3 unanchored, sampled-trained, seed 123 & 50 & 30.7 & 9.2 & 116.2 & 0.271 & 0.756 & 135 / 0 & 0.29 \\
greedy: B3 unanchored, sampled-trained, pooled & 100 & 18.3 & 9.5 & 357.6 & 0.330 & 0.743 & 252 / 0 & 0.23 \\
greedy: B4 unanchored, greedy-trained, seed 42 & 50 & 6.0 & 9.8 & 619.5 & 0.370 & 0.704 & 2306 / 0 & 0.17 \\
greedy: B4 unanchored, greedy-trained, seed 123 & 50 & 23.3 & 9.4 & 161.8 & 0.461 & 0.931 & 9486 / 0 & 0.22 \\
greedy: B4 unanchored, greedy-trained, pooled & 100 & 14.7 & 9.6 & 409.3 & 0.415 & 0.817 & 11792 / 0 & 0.20 \\
\bottomrule
\end{tabular}}
\vspace{1mm}
\parbox{\textwidth}{\scriptsize Arms B1/B2 are anchored (CALM), B3/B4 unanchored, all from the identical prior checkpoint. The sampled rows are the primary comparison (the rule CALM trains with); under greedy execution B1 and B2 are the constant Lane\_Right policy.}
\end{table*}

\begin{table*}[t]
\centering
\caption{The pre-registered decision rule (protocol \S8.9 rev.~2) applied to B1 vs B3 under sampled execution and mode 512: primary $\Delta H_\mathrm{sim} \geq 0.10$ and $\Delta H_\mathrm{val} \geq 0.05$ on every seed; collapse $H_\mathrm{sim} \leq 0.50$ or $H_\mathrm{val} \leq 0.60$; guards by paired bootstrap over matched episodes within seed (10,000 resamples): exits difference lower bound $\geq -0.25$, proximity ratio upper bound $\leq 1.25$. Outcome: pass.}
\label{tab:t5b_decision}
\footnotesize
\resizebox{\textwidth}{!}{%
\begin{tabular}{lcccccccc}
\toprule
\textbf{Seed} & \textbf{$H_\mathrm{sim}$ B1} & \textbf{$H_\mathrm{sim}$ B3} & \textbf{$\Delta H_\mathrm{sim}$} & \textbf{$H_\mathrm{val}$ B1} & \textbf{$H_\mathrm{val}$ B3} & \textbf{Exits/ep B1 vs B3, 95\,\% CI of diff.} & \textbf{Prox./km B1 vs B3, 95\,\% CI of ratio} & \textbf{Primary / exits guard / prox. guard} \\
\midrule
seed 42 & 0.714 & 0.254 & 0.460 & 0.809 & 0.632 & 2.56 vs 1.16, [+1.12, +1.68] & 37.4 vs 482.9, [0.061, 0.100] & pass / pass / pass \\
seed 123 & 0.706 & 0.216 & 0.490 & 0.806 & 0.450 & 2.40 vs 1.36, [+0.70, +1.34] & 41.0 vs 121.9, [0.290, 0.385] & pass / pass / pass \\
\bottomrule
\end{tabular}}
\end{table*}


\subsection{Against the frozen prior}
\label{sec:prior}
\todo{Protocol \S8.11 around Table~\ref{tab:t6_frozen_prior}; Fig.~F4 paired bootstrap distributions.}
\begin{table*}[t]
\centering
\caption{Anchored CALM against the frozen prior it was initialised from, on matched episodes (protocol \S8.11): the Stage 0 class-balanced prior loaded into all five agents and never fine-tuned, evaluated with the Option B episode seeds under mode 512. Differences are CALM minus prior with 95\,\% paired-bootstrap intervals over episodes within seed (10,000 resamples); proximity as the ratio CALM/prior of pooled rates.}
\label{tab:t6_frozen_prior}
\footnotesize
\resizebox{\textwidth}{!}{%
\begin{tabular}{lcccccc}
\toprule
\textbf{Execution: arm vs prior} & \textbf{$n$} & \textbf{$H$ per-state sim} & \textbf{Exits/ep} & \textbf{Speed m/s} & \textbf{Prox./km ratio} & \textbf{Lane changes/ep} \\
\midrule
sampled: B1 vs prior, seed 42 & 50 & 0.734 $\to$ 0.713, [-0.028, -0.014] & 2.98 $\to$ 2.56, [-0.62, -0.24] & 11.7 $\to$ 11.9, [-1.0, +1.4] & $\times$0.770 [0.631, 0.938] & 359 $\to$ 10, [-372, -327] \\
sampled: B1 vs prior, seed 123 & 50 & 0.728 $\to$ 0.705, [-0.031, -0.015] & 2.98 $\to$ 2.40, [-0.82, -0.34] & 12.2 $\to$ 11.2, [-1.9, +0.1] & $\times$0.905 [0.764, 1.074] & 324 $\to$ 10, [-337, -293] \\
sampled: B2 vs prior, seed 42 & 50 & 0.734 $\to$ 0.720$^\dagger$ & 2.98 $\to$ 2.56, [-0.62, -0.24] & 11.7 $\to$ 11.8, [-1.0, +1.2] & $\times$0.787 [0.647, 0.959] & 359 $\to$ 10$^\dagger$ \\
sampled: B2 vs prior, seed 123 & 50 & 0.728 $\to$ 0.707$^\dagger$ & 2.98 $\to$ 2.08, [-1.18, -0.64] & 12.2 $\to$ 10.4, [-2.7, -0.9] & $\times$0.984 [0.856, 1.135] & 324 $\to$ 10$^\dagger$ \\
\midrule
greedy: B1 vs prior, seed 42 & 50 & 0.743 $\to$ 0.649, [-0.102, -0.085] & 2.08 $\to$ 3.00, [+0.84, +0.98] & 17.4 $\to$ 24.5, [+6.3, +7.8] & $\times$0.125 [0.110, 0.143] & 259 $\to$ 10, [-269, -232] \\
greedy: B1 vs prior, seed 123 & 50 & 0.739 $\to$ 0.636, [-0.114, -0.091] & 2.08 $\to$ 3.00, [+0.84, +0.98] & 17.3 $\to$ 24.5, [+6.4, +8.1] & $\times$0.120 [0.103, 0.143] & 248 $\to$ 10, [-254, -223] \\
greedy: B2 vs prior, seed 42 & 50 & 0.743 $\to$ 0.666$^\dagger$ & 2.08 $\to$ 3.00, [+0.84, +0.98] & 17.4 $\to$ 24.5, [+6.3, +7.8] & $\times$0.125 [0.110, 0.143] & 259 $\to$ 10$^\dagger$ \\
greedy: B2 vs prior, seed 123 & 50 & 0.739 $\to$ 0.658$^\dagger$ & 2.08 $\to$ 3.00, [+0.84, +0.98] & 17.3 $\to$ 24.5, [+6.4, +8.0] & $\times$0.120 [0.104, 0.142] & 248 $\to$ 10$^\dagger$ \\
\bottomrule
\end{tabular}}
\vspace{1mm}
\parbox{\textwidth}{\scriptsize $^\dagger$B2's per-episode step log was overwritten by a later arm (protocol \S8.10 caveat); its entropy and lane changes are the per-seed summary values without an interval. B1 was re-evaluated with its own log, identical to the 4 November run. Under sampled execution the prior has equal-or-higher per-state entropy, more exits and the same speed; proximity favours CALM on seed 42 only. Under greedy execution B1 and B2 are the constant Lane\_Right policy (Table~\ref{tab:t3_gate}).}
\end{table*}


\subsection{Training reward, labelled}
\label{sec:reward}
\todo{Fig.~F5 training-reward curves for the four Option B arms, stated as a training statistic the evaluator does not measure.}

\section{What This Benchmark Can and Cannot Test}
\label{sec:testability}
\todo{Table T7 = review \S4, with the scenario facts: only lane 0 of highway\_seg1 reaches the ramp; two of five agents are passive mainline vehicles; NGSIM I-80 is a congested bottleneck at about 6 m/s, MultiExit free-flow.}

\section{Discussion}
\label{sec:discussion}
\todo{Regulariser that preserves, not improves; the circularity concern and what remains non-trivial; reporting checklist; limitations: two seeds for the isolation, one scenario, no entropy sweep or MASAC on the corrected contract, reward training-time only, the prior's class weighting determines what is preserved.}

\section{Conclusion}
\label{sec:conclusion}
\todo{The three measurable findings; a benchmark in which the exit policy is non-trivial as the next step.}

\appendix
\section*{Historical record}
\todo{Decision 5 on \#67 pending: the submitted Tables I, III, IV, V as published, labelled as measurements of the defective pipeline, with the defect responsible for each column.}

\bibliographystyle{IEEEtran}
\bibliography{../references}

\end{document}
